\section{Baselines and Dual-Solver Details}
\label{app:baselines}
%==============================================================================

\subsection{Dual calibration}
\label{ssec:dual-fix}
The solver works entirely in log-space: $p$ is never materialized as a probability tensor; $q_\lambda$'s logits are $\log p-\lambda\risk$ computed from \texttt{F.log\_\allowbreak{}softmax} of the temperature-scaled logits, and $\KL(q\Vert p)$ and both entropies are computed from log-probabilities, with $0\log0:=0$ and a conversion to bits (division by $\ln2$). This matters at low temperature: at $T{=}0.10$ many vocabulary entries have probabilities below the float32 range, and materializing $p$ would collapse distinct low-probability tokens to a single floor value.

Structural infeasibility is decided in closed form before any search: if $\min_v\risk(v)>\varepsilon$ (with tolerance), the step is reported as \texttt{structurally\_\allowbreak{}infeasible} and no search is attempted. Otherwise a feasible $\lambda$ exists, and the search bracket expands geometrically (doubling from $\lambda{=}1$, at most 60 doublings, so $\lambda\le2^{60}\approx1.15\times10^{18}$) until its upper end is feasible, and then 30 bisection iterations are run within it. The limit of 60 doublings is practical and not mathematical; Section~\ref{ssec:lambda-estimate} shows that it is far beyond anything a risk-weighted logit shift requires away from the boundary case.

\textbf{Numerical tolerances.} Risk values are float32 and $\varepsilon$ is float64, so the comparison of $\min_v\risk(v)$ with $\varepsilon$ uses an absolute tolerance of $10^{-6}$ (in risk units), and a step with $|\min_v\risk(v)-\varepsilon|\le10^{-6}$ is flagged \texttt{near\_\allowbreak{}boundary}, which is distinct from both an interior point and structural infeasibility; in this case the feasibility target is relaxed by the same tolerance. Per-step diagnostics (feasibility status, $\lambda$, achieved risk, violation, final bracket width, number of doublings, and the boundary flag) are saved for every dual-mode configuration. On the primary model the dual solver was feasible at 100.0\% of steps for both $\varepsilon$, with zero mean and zero maximum violation at the stored precision of four decimals, a mean of about 2.2 bracket doublings per step and a maximum of 9, and no near-boundary or structurally infeasible step. In the exact systems the one-step solver (\texttt{one\_\allowbreak{}step\_\allowbreak{}projection.py}) expands the bracket for $\lambda$ geometrically from 1 until the budget is met (capped at $10^{9}$, beyond which it raises an error) and then bisects for 100 iterations. It classifies each step as inactive, interior, boundary or infeasible with a risk tolerance of $10^{-12}$, applies the rule of Section~\ref{sec:setup} to the last two, checks that the final distribution meets the budget, and raises an error if it does not, so a failed solve is never counted as a success (\texttt{tests/test\_\allowbreak{}infeasible\_\allowbreak{}rule.py}). The oracle and the exact forward propagations carry residuals below $10^{-10}$.

\subsection{Recurrence risk and prompt context}
$\risk(v)$ uses an incrementally maintained hash map, checked to be identical to brute-force history scanning. Whether the prompt's own tokens count as history was tested (\texttt{include\_\allowbreak{}prompt\_\allowbreak{}context}$\in\{$True,False$\}$): the RMST difference (492.9 vs.\ 499.9) is not significant under the primary exact test ($p{=}0.125$).

\subsection{The adaptive controller}
$\lambda_t=\lambda_{\text{base}}+\lambda_{\text{rep}}\max(0,\overline{\text{rep}}-\text{rep}_{\text{target}})-\lambda_{\text{ent}}\max(0,\overline{H}-H_{\text{target}})$, clipped to $[0,\lambda_{\max}]$. The lower clip is a correctness requirement, since an unclipped negative multiplier would reward risk.

\subsection{SuffixMatch}
\textsc{SuffixMatch} finds the longest suffix of the generated history (last 400 tokens) that occurred earlier and subtracts $\alpha\cdot\text{match\_len}/20$ from the logits of the tokens that followed those earlier occurrences, with $\alpha{=}3$; it scans the history at every step, at cost $O(\text{history})$. It is a homemade longest-suffix-match heuristic inspired by LZ77, not the published Look-back decoding algorithm and not attributed to any publication.

\subsection{FSD-style decoder}
\label{app:fsd}
Frustratingly Simple Decoding \citep{yang2023frustratingly} contrasts the model with a training-free $n$-gram predictor built from the text generated so far (the anti-LM). The authors' code is at github.com/LHRYANG/FSD; the version used for comparison is commit \texttt{46df7d5bc622\allowbreak{}8bd5045b68c0\allowbreak{}9f85d75bcceb\allowbreak{}1fdf} (\texttt{ngram\_model/fsd.py} last changed in \texttt{0f469ac54f26\allowbreak{}b02d20f831a8\allowbreak{}27c09336ede5\allowbreak{}f06a}, 2024-02-23). At that commit \texttt{utils.py} calls its \texttt{NGram} objects directly, which the same commit's \texttt{fsd.py} does not define, so the penalty is checked against \texttt{NGram.penalize}, the function that computes it.

The FSD-style decoder here follows the official scoring. Candidates are the top $k{=}6$ tokens by $p_\theta$ and are ranked by $p_\theta(v)-\alpha\,p_\omega(v)$ with $\alpha{=}3$ (the official code writes $(1{-}a)p_\theta-a\,p_\omega$ with $a=\alpha/(1{+}\alpha)$, which has the same argmax). The anti-LM $p_\omega$ is the recursive residual-mass construction: for a candidate, orders are visited from the highest ($N{=}3$, context length 2) down to the unigram, with $\text{remaining}=1$; an order at which the candidate never followed the current context is skipped and leaves its residual to the lower orders; an order at which it did contributes $\text{remaining}\cdot\beta\cdot\text{count}/(\text{total}{+}1)$ ($\beta{=}0.9$), and the unigram contributes $\text{remaining}\cdot\text{count}/\text{total}$ without $\beta$; after each contributing order $\text{remaining}\leftarrow\text{remaining}(1-\beta)$. When every order contributes the weights are about $0.90/0.09/0.01$. The anti-LM scores of this decoder equal those of \texttt{NGram.penalize} to within $10^{-9}$ on constructed prefixes and 300 random histories (\texttt{tests/test\_\allowbreak{}fsd\_\allowbreak{}against\_\allowbreak{}official.py}, which fetches the pinned file at run time because the repository carries no license).

It departs from the official method in two respects, so it is labeled FSD-style (adapted) and is not presented as a reproduction. The official method takes the argmax over the top-$k$ scores, whereas this study resamples every configuration over several seeds, so the decoder samples from the top-$k$ scores after clamping at zero and renormalizing (from $p_\theta$'s own top-$k$ if every score clamps to zero). And $p_\theta$ is the temperature-scaled softmax that serves as the reference distribution throughout, where the official code uses the raw softmax.

\subsection{LZ Penalty}
\label{app:lz}
The LZ Penalty \citep{ginart2025lzpenalty} penalizes each candidate token by the change in codelength it would induce under an LZ-style compressor applied to a recent buffer against a lookback window. The reference implementation is github.com/tginart/sglang, \texttt{python/\allowbreak{}sglang/\allowbreak{}srt/\allowbreak{}sampling/\allowbreak{}penaltylib/\allowbreak{}lz\_\allowbreak{}penalty.py} at commit \texttt{ab08e6976d49\allowbreak{}80e350df3891\allowbreak{}e40967a9718f\allowbreak{}ae8b} (2025-07-01, Apache-2.0). For a context $W$ and buffer length $n$, the penalty vector of its \texttt{compute\_lz\_penalty\_stateless} is $\log_2V-\min\{\max\{\log_2 d(a),0\},\log_2V\}$ for every candidate $a$, where $d(a)$ is the distance to the closest position of the context at which $a$ follows the longest match of the buffer's trailing tokens ($+\infty$, hence penalty $0$, if $a$ never occurred), and the reference subtracts it from the logits scaled by the user weight. This study adds $\alpha\cdot(\text{cost}-\log_2V)$ to the logits, the same quantity with the opposite sign.

\texttt{LZPenaltyDecoder.\_penalty\_vector} equals the reference function exactly (maximum difference $0.0$) on constructed and 300 random contexts (\texttt{tests/test\_\allowbreak{}lz\_\allowbreak{}against\_\allowbreak{}reference.py}). Two compatibility notes apply to that test. The reference casts an integer tensor with \texttt{torch.full\_like(x, inf)}, which raises on current PyTorch, so the extracted function is run with an added float cast. And its early return for a context no longer than $n$ passes a tensor as a size to \texttt{torch.zeros} and raises, so that branch is checked on this side only (an all-zero penalty).

The decoding path differs from the reference's state handling, so this decoder is labeled LZ Penalty (adapted): the reference keeps only generated tokens in a window of \texttt{lz\_\allowbreak{}lookback\_\allowbreak{}size}, whereas this implementation keeps the prompt and the generated tokens and uses the last \texttt{buffer\_\allowbreak{}size}$+$\texttt{window\_\allowbreak{}size} of them. Buffer and window are 8 and 128 with $\alpha{=}0.15$ for the character model, and 32 and 512 for the subword models, scaled down for the character vocabulary, and the penalty is applied to the logits before temperature scaling.

\subsection{Order-of-magnitude estimate for the multiplier cap}
\label{ssec:lambda-estimate}
\label{app:lambda-estimate}
For a two-outcome simplification (one risky token with $\risk{=}1$, one safe token with $\risk{=}0$, raw pre-temperature probability ratio $\rho$ risky:safe), the $\lambda$ needed to bring $q(\text{risky})\le\varepsilon$ at temperature $T$ is $\lambda\ge\ln\frac{1-\varepsilon}{\varepsilon}+\frac{1}{T}\ln\rho$. At $T{=}0.10$, $\rho{=}100$ needs $\lambda\approx50$ and $\rho{=}1000$ needs $\lambda\approx73$. This is for an interior point, away from the $\varepsilon=\min_v\risk(v)$ boundary, where the required $\lambda$ depends on how small a float64 gap the solver must resolve against a float32-rounded risk floor and is handled by the \texttt{near\_\allowbreak{}boundary} tolerance. The 60-doubling limit exceeds the $\rho{=}1000$ estimate by roughly 16 orders of magnitude and is not observed to bind outside the boundary case.

